\ifdefined\ishandout
\documentclass[handout]{beamer}
\else
\documentclass{beamer}
\fi

\usepackage[frenchb]{babel}
\usepackage[T1]{fontenc}
\usepackage[latin1]{inputenc}
\usepackage{hyperref}
\usepackage{multirow}
\usepackage{listings}
\usepackage{fancyvrb}
\usepackage{tikz}
\usepackage{framed}
\usepackage{algorithm}
\usepackage{algorithmic}
\usepackage{xcolor}
\usepackage{color, colortbl}
\usepackage{handoutWithNotes}
\usepackage{amsmath}
\usetikzlibrary{shapes.geometric}
\usetikzlibrary{positioning}
\usetikzlibrary{shapes.arrows, chains}
\usetikzlibrary{arrows,calc}
\usetikzlibrary{shapes.multipart}
\usepackage{array}
\usetheme{Boadilla}
\definecolor{BlueGreen}{cmyk}{0.85,0,0.33,0}
\definecolor{Gray}{rgb}{0.8,0.8,0.8}

\ifdefined\ishandout
\pgfpagesuselayout{3 on 1 with notes}[a4paper,border shrink=5mm]
\usecolortheme{dove}
\else
\usecolortheme{dolphin}
\fi


\lstnewenvironment{codeC}
{ \lstset{language=C,
    otherkeywords={printf,scanf}}
}
{}

\ifdefined\ishandout
\definecolor{mygreen}{rgb}{0,0,0}
\definecolor{mymauve}{rgb}{0,0,0}
\definecolor{myblue}{rgb}{0,0,0}
\else
\definecolor{mygreen}{rgb}{0,0.6,0}
\definecolor{mymauve}{rgb}{0.58,0,0.82}
\definecolor{myblue}{rgb}{0,0,1}

\fi

\definecolor{mygray}{rgb}{0.5,0.5,0.5}


\lstset{language=C,
% breakatwhitespace=false,         % sets if automatic breaks should only happen at whitespace
%  breaklines=true,                 % sets automatic line breaking
%  captionpos=b,                
commentstyle=\itshape\color{mymauve},
keywordstyle=\bfseries\color{myblue},
%numbers=left,                    % where to put the line-numbers; possible values are (none, left, right)
%  numbersep=8pt,                   % how far the line-numbers are from the code
%  numberstyle=\tiny\color{mygray}, % the style that is used for the line-numbers
  rulecolor=\color{black},         % if not set, the frame-color may be changed on line-breaks within not-black text (e.g. comments (green here))
%  showspaces=false,                % show spaces everywhere adding particular underscores; it overrides 'showstringspaces'
  showstringspaces=false,          % underline spaces within strings only
%  showtabs=false,                  % show tabs within strings adding particular underscores
%  stepnumber=2,                    % the step between two line-numbers. If it's 1, each line will be numbered
  stringstyle=\color{mygreen},     % string literal style
%  tabsize=2 
}
\ifdefined\ishandout
\newcommand{\red}{\textbf}
\else
\newcommand{\red}{\textcolor{red}}
\fi
%\newcommand \emph
%Default size : 12.8 cm * 9.6 cm

\newcommand{\tmark}[1]{\tikz[remember picture, baseline=-.5ex]{\coordinate(#1);}}

\ifdefined\ishandout
\newenvironment<>{codeblock}[1]{%begin
  \setbeamercolor{block title}{fg=black,bg=lightgray!80}%
  \begin{block}{#1}}
  % \begin{codeC}}
  %  {\end{codeC}
{  
\end{block}}

\newenvironment<>{termblock}[1]{
    \setbeamercolor{block title}{fg=black,bg=lightgray!90}%
    \begin{block}{#1}
}
%     \begin{Verbatim}}
{%\end{Verbatim}
\end{block}
}

\definecolor{bluegreen}{RGB}{0,0,0}
%\definecolor{bluegreen}{rgb}{0,0.6,0.8}
\else

\newenvironment<>{codeblock}[1]{%begin
  \setbeamercolor{block title}{fg=darkgray,bg=yellow}%
  \begin{block}{#1}}
  % \begin{codeC}}
  %  {\end{codeC}
{  
\end{block}}

\newenvironment<>{termblock}[1]{
    \setbeamercolor{block title}{fg=white,bg=lightgray}%
    \begin{block}{#1}}
%     \begin{Verbatim}}
{%\end{Verbatim}
\end{block}
}

\definecolor{bluegreen}{RGB}{0,149,182}
%\definecolor{bluegreen}{rgb}{0,0.6,0.8}
\fi

%\newcommand{\output}[1]{
\setbeamertemplate{navigation symbols}{}
\newcommand{\bvrb}{\Verb[commandchars=���,formatcom=\color{bluegreen}]}
\newcommand{\footvrb}{\footnotesize\Verb}
\newcommand{\vrbalert}[2][]{\visible<#1>{#2}}
%%% Commande pour les listes/arbres
\newcommand{\mvide}{\nodepart{one} \nodepart{two}}
\newcommand{\tvide}{\nodepart{one} \nodepart{two} \nodepart{three}}

%%Fin des commandes pour les listes/arbres.


\newcommand<>{\case}[2]{%
\filldraw (#1,#2) rectangle (1+#1,1+#2);}
%%% Param�tres du cours (� r�gler)
%Num�ro du cours
\newcommand{\nb}{0}
\title[Informatique]{Pr�sentation du cours d'informatique MAIN3/ROB3}
\author[]{julien.brajard@upmc.fr}
\institute[Polytech'UPMC]{Polytech'UPMC}
\date{18 Septembre 2017}
\begin{document}

\begin{frame}
\titlepage
\centering{
\url{https://moodle-sciences.upmc.fr} (cours Informatique G�n�rale MAIN-ROB)}
\end{frame}

%%%%%%%%%%%%%%%%%%%%%
\begin{frame}
\frametitle{Objectif du cours}
\begin{alertblock}{}
Le cours d'informatique g�n�ral a pour objectif de vous 
initier � l'algorithmique et aux concepts de base de 
la programmation � travers le langage C.
\end{alertblock}

\end{frame}

%%%%%%%%%%%%%%%%%%%%%%%%%%%%%%%%%%%%%%%%%%%
\begin{frame}
\frametitle{D�roulement du cours}

\begin{itemize}
\item 10 cours de 2h,\\
\item 6 TD de 2h,\\
\item 30H TP\\
\end{itemize}
\end{frame}

\begin{frame}
\frametitle{Planning des ROB}
\vspace{-2em}
\begin{table}[h]
\centering
\begin{tabular}{|l|l|l|l|}
\hline
\rowcolor{BlueGreen}
Semaine & Cours & TP & TD \\
\hline
S38 - 18/09 & CM1  & \free & TD1 \\
\hline
S39 - 25/09 & CM2 (fac.) & \free  & \free \\
\hline
S40 - 02/10 & CM3 + \textbf{CC1}   & TP1  & TD2  \\
\hline
S41 - 09/10 & CM4  & \free  & TD3  \\
\hline
S42 - 16/10 & CM5 & TP2 & TD4 \\
\hline
S43 - 23/10 & CM6  & TP3  & TD5  \\
\hline
%S44 - 31/10 & \free & \free & \free \\
%\hline
S45 - 06/11 & CM7   & TP4   & \\
\hline
S46 - 13/11 & CM8  & TP5 & \\
\hline
S47 - 20/11 & CM9 & TP6 & \textbf{Soutenance orale} \\
\hline
S48 - 27/11 & \textbf{CC2}  & Pr. 1 &\free  \\
\hline
S49 - 04/12 & \free & Pr. 2 & \free \\
\hline
 S50 - 11/12 & \free & Pr. 3 & \free \\
\hline
 S51 - 18/12 & \free & Pr. 4 & \free \\
\hline
S02 - 08/01 & \free &\textbf{Soutenance Projet}  & \free \\
\hline
S?? - ??/?? & \multicolumn{3}{|c|}{Test du projet (facultatif)} \\
\hline
\end{tabular}
\end{table}


\end{frame}

\begin{frame}
\frametitle{Planning des MAIN}
\vspace{-2em}
\begin{table}[h]
\centering
\begin{tabular}{|l|l|l|l|}
\hline
\rowcolor{BlueGreen}
Semaine & Cours & TP & TD \\
\hline
S38 - 18/09 & CM1  & \free & TD1 \\
\hline
S39 - 25/09 & CM2 (fac.) &\free  & \free \\
\hline
S40 - 02/10 & CM3 + \textbf{CC1}   & TP1  & TD2  \\
\hline
S41 - 09/10 & CM4  & TP2  & TD3  \\
\hline
S42 - 16/10 & CM5 & TP3 & TD4 \\
\hline
S43 - 23/10 & CM6  & TP4  & TD5  \\
\hline
%S44 - 31/10 & \free & \free & \free \\
%\hline
S45 - 06/11 & CM7   & TP5   & \textbf{Soutenance orale}\\
\hline
S46 - 13/11 & CM8  & \free& \\
\hline
S47 - 20/11 & CM9 & Pr. 1 & \\
\hline
S48 - 27/11 & \textbf{CC2}  & \free &\free  \\
\hline
S49 - 04/12 & \free & Pr. 2 & \free \\
\hline
 S50 - 11/12 & \free & \free & \free \\
\hline
 S51 - 18/12 & \free & Pr. 3 & \free \\
\hline
S02 - 08/01 & \free &\textbf{Soutenance Projet}  & \free \\
\hline
S?? - ??/?? & \multicolumn{3}{|c|}{Test du projet (facultatif)} \\
\hline
\end{tabular}
\end{table}


\end{frame}


\begin{frame}
\frametitle{Evaluation}
\begin{itemize}

\item QCM en ligne � la fin des cours.\\
\begin{itemize}
\item � remplir en ligne sur Moodle dans les deux semaines qui suivent le cours.\\
\end{itemize}

\item Contr�le continu sur papier\\
\begin{itemize}
\item Contr�le de 30 min le 02/10 d'�valuation des acquis\\
\item Contr�le d'1h45 le 27/11\\
\end{itemize}

\item Travail "maison"\\
\begin{itemize}
\item Travail d'une dur�e de 1 heure � faire entre le 09/10 et le 06/11.\\
\end{itemize}

\item Soutenance en groupes\\
\begin{itemize}
\item Pr�sent� en TD la semaine du 23 Octobre. \\
\item Soutenance en TD la semaine du 6 Novembre (MAIN) ou 20 Novembre (ROB).\\
\end{itemize}

\item Projet de fin de semestre\\
\begin{itemize}
\item Une soutenance orale obligatoire la semaine du  08 Janvier 2018\\
\item Un test du projet avec les volontaires (date � d�terminer)\\
\end{itemize}
\end{itemize}


\end{frame}

\begin{frame}
\frametitle{Plateforme d'enseignement en ligne Moodle}

\begin{figure}
\raggedright
\hfill \includegraphics[scale=0.1]{./fig/moodle.jpg}
\end{figure}

\begin{itemize}
\item Accessible sur~\url{moodle-sciences.upmc.fr}.\\
\item Connection avec votre login et mot de passe annuaire.\\
\item Onglet \textbf{Informatique G�n�rale MAIN-ROB}.\\
\end{itemize}
\begin{alertblock}{}
Vous devez pouvoir acc�der � votre e-mail dans l'annuaire de l'UPMC.
\end{alertblock}
\end{frame}

\end{document}